\section*{第 \ref{ch:b3:learning-memory} 章——神经可塑性、学习与记忆}

\begin{solution}{exo:b3:learning-memory:1}
\omterm{def:b3:learning-memory:kinds}{陈述性记忆}（事实、事件）在 H.M. 身上丧失——他
无法回忆任何新近内容，而非\omterm{def:b3:learning-memory:kinds}{陈述性记忆}则完好保留：他
能正常学会镜像描画，却记不得练习过程。技能：
依赖纹状体和\omterm{def:b3:nervous-systems:plan}{小脑}；条件性恐惧：杏仁核；条件性运动
反射：\omterm{def:b3:nervous-systems:plan}{小脑}；启动效应：感觉皮层。
\end{solution}

\begin{solution}{exo:b3:learning-memory:2}
当突触前细胞反复参与使突触后细胞放电时，
二者之间的连接即得到增强。因此\omterm{def:b3:learning-memory:ltp}{长时程增强}具有输入特异性（只有
参与其中的活动突触发生改变），属联合型（弱输入在
强输入使细胞放电的同时处于活动，也会得到增强），
并且具有\omterm{def:b3:structural-biology:allostery}{协同性}（需要足够多的输入共同作用才能使细胞放电）。
\end{solution}

\begin{solution}{exo:b3:learning-memory:3}
该受体仅在谷氨酸结合（突触前细胞已放电）
\emph{并且}膜去极化到足以排出镁离子（突触后细胞处于活动）时才导通：
赫布法则的两个条件集于同一蛋白质。
高频刺激前施加 AP5：无钙内流，无\omterm{def:b3:learning-memory:ltp}{长时程增强}，
尽管经 AMPA 受体的传递正常。高频刺激后施加 AP5：
已诱导的\omterm{def:b3:learning-memory:ltp}{长时程增强}不受影响——该受体是诱导所需，
而非维持所需。
\end{solution}

\begin{solution}{exo:b3:learning-memory:4}
短时程：对已有蛋白质进行共价修饰——\omterm{prop:b3:learning-memory:aplysia}{PKA}
磷酸化并关闭一个钾通道，使动作电位加宽，
并增加递质释放；不需要新蛋白质。长时程：\omterm{prop:b3:learning-memory:aplysia}{PKA} 激活 \omterm{prop:b3:learning-memory:nmda}{CREB}，
新基因被转录，新的突触末端生长。在反复训练期间给予
蛋白质合成抑制剂，短时程的易化效应保持不变，
但会消除数天后测得的记忆；
若在训练之后给予，则毫无作用。
\end{solution}

\begin{solution}{exo:b3:learning-memory:5}
$0.7^{n} \le 0.1$：$n \ge 6.5$，故需 $7$ 次试验。从 $0.9\lambda$ 出发，
$0.9\times 0.7^{n} < 0.1$：同样需 $7$ 次试验。两者对称，因为
同一 $\alpha$ 支配两个方向。在动物中，消退通常
较慢，且原有记忆在休息之后会重新出现（自发恢复）：
消退是一种新的学习，它抑制旧记忆，而非将其
抹除。
\end{solution}

\begin{solution}{exo:b3:learning-memory:6}
特征向量为 $(1,1)/\sqrt{2}$（本征值 $1.5$）和 $(1,-1)/\sqrt{2}$，
对应本征值 $0.5$；赫布法则使 $\mathbf{w}$ 转向 $(1,1)$，神经元
由此能够检测两路输入同时放电。奥贾更新的一步：$y = 1$，
$\Delta\mathbf{w} = 0.1\times 1\times((1,1) - 1\times(1,0)) = (0,0.1)$，
因此 $\mathbf{w} = (1, 0.1)$。
\end{solution}

\begin{solution}{exo:b3:learning-memory:7}
$500\times 0.05 = 25$ 个游离离子；$25/6\times 10^{23} = 4.2\times 10^{-23}$
摩尔，溶于 $10^{-16}$ L 中：\qty{0.42}{\micro mol/L}，为静息水平的四倍，
但低于钙调蛋白的 $K_{d}$——仅部分激活；需要
若干通道同时开放，才能达到
使 \omterm{prop:b3:learning-memory:nmda}{CaMKII} 开启的微摩尔浓度范围。
\end{solution}

\begin{solution}{exo:b3:learning-memory:8}
训练 A 直至 $V_{A} = \lambda$：早期试验中\omterm{prop:b3:learning-memory:rescorla}{多巴胺}爆发，
随着奖励变得可预测，\omterm{prop:b3:learning-memory:rescorla}{多巴胺}转为沉默。对复合
刺激试验而言，预测 $V_{A} + V_{B} = \lambda$ 已经精确，
误差为零，\omterm{prop:b3:learning-memory:rescorla}{多巴胺}沉默，$V_{B}$ 从不改变：B 被
阻断。单独检验时，B 不引起任何反应；若此后单独给予 B 一个奖励，
则该奖励是意料之外的，\omterm{prop:b3:learning-memory:rescorla}{多巴胺}爆发。
\end{solution}

\begin{solution}{exo:b3:learning-memory:9}
(a) \omterm{def:b3:learning-memory:hippocampus}{CA1} 区无\omterm{def:b3:learning-memory:ltp}{长时程增强}，并出现空间记忆缺陷，其他学习
能力正常。(b) 短时记忆正常，长时记忆在
\qty{24}{h} 时缺失。(c) 无影响——\omterm{def:b3:learning-memory:kinds}{巩固}已完成（除非记忆
被再激活，此时它又变得不稳定）。(d) 小鼠在安全环境中
出现僵住反应：该记忆是通过再激活
编码它的细胞而被人为唤起的。
\end{solution}

\begin{solution}{exo:b3:learning-memory:10}
$\Delta|\mathbf{w}|^{2} \approx 2\mathbf{w}\cdot\Delta\mathbf{w} = 2\eta
y^{2} \ge 0$：长度永不减小，且每当神经元
放电时都会增大。在奥贾不动点处 $C\mathbf{w} = (\mathbf{w}^{\mathsf{T}}C
\mathbf{w})\mathbf{w}$，因此 $\mathbf{w}$ 是特征向量，$\lambda =
\lambda|\mathbf{w}|^{2}$，故 $|\mathbf{w}| = 1$。无界增长
会使每个突触饱和、消除选择性，并引发
失控兴奋；生物学中的归一化机制包括突触缩放（神经元
的所有突触在其过度活动时按比例下调）、\omterm{def:b3:learning-memory:ltp}{长时程抑制}，以及
睡眠期间的突触下调。
\end{solution}

\begin{solution}{exo:b3:learning-memory:11}
大鼠：$0.14\times 3\times 10^{5} \approx 4\times 10^{4}$；人：$0.14
\times 2\times 10^{6} \approx 3\times 10^{5}$。该估计假设
模式随机且稠密、连接完全；实际模式是稀疏的
（这会大大提高容量），且具有结构性，而\omterm{def:b3:learning-memory:kinds}{海马}是
一个暂时性存储，它将记忆转交给由 $10^{10}$ 个神经元组成的皮层
——终身存储位于皮层，\omterm{def:b3:learning-memory:kinds}{海马}中的数值只是
缓冲区大小，而非记忆数量。
\end{solution}

\begin{solution}{exo:b3:learning-memory:12}
睁开眼的输入与皮层细胞的放电相关，
因而得到增强；闭合眼的输入则不然，在
总权重有限的竞争中，它们会减弱——即赫布法则加上
归一化。双眼闭合时，两路输入均未受青睐，因此
两者都保留（较弱的）领地，效应较为温和。\omterm{def:b3:learning-memory:critical}{关键期}
随抑制性回路成熟而结束，神经周网
在突触周围变硬，NMDA 受体的亚基转变为
可塑性较低的形式。
\end{solution}

\begin{solution}{pb:b3:learning-memory:1}
\textbf{1.} $10\times 10^{3} = 10^{4}$ 个离子；\qty{2}{\%} 游离，得 $200$；
$200/6\times 10^{23} = 3.3\times 10^{-22}$ 摩尔，位于 $8\times 10^{-17}$ 升中：
\qty{4.2}{\micro mol/L}。
\textbf{2.} 是的，高于 \qty{2}{\micro mol/L} 的阈值。两
个受体：$40$ 个游离离子，\qty{0.83}{\micro mol/L}——低于阈值。
\textbf{3.} 在 \qty{0.83}{\micro mol/L} 时，钙调神经磷酸酶活跃，而 \omterm{prop:b3:learning-memory:nmda}{CaMKII}
不活跃：AMPA 受体被移除，突触减弱（\omterm{def:b3:learning-memory:ltp}{长时程抑制}，LTD）。一
个离子，两种亲和力不同的酶：较大而快速的上升抵达
激酶，较小而缓慢的上升则只抵达磷酸酶。
\textbf{4.} 在数个时间常数之内到达的第二个输入，会将其钙
加到第一个输入残留的钙上；经过 \qty{50}{ms} 后，几乎
没有残留。NMDA 的同时性要求与 \qty{20}{ms} 的衰减
共同给出两侧各约 \qty{20}{ms} 的窗口——即
峰电位时序窗口。
\textbf{5.} 缓冲液在不到
一微米的范围内捕获大部分钙，狭窄的棘颈减慢了钙向树突的扩散，
因此钙的升高局限于活跃的那一个棘突；
相距 \qty{0.5}{\micro m} 的邻近棘突毫无感知，
只有活跃的突触发生改变。
\textbf{6.} $(70 - 30)/5 = 8$ 路同步输入，或来自细胞体的逆向传播
动作电位：单个突触无法解除
其自身 NMDA 受体的阻断——\omterm{def:b3:structural-biology:allostery}{协同性}由其物理性质决定。
\textbf{7.} $(1,1)/\sqrt{2}$ 的本征值为 $1.6$；$(1,-1)/\sqrt{2}$
的本征值为 $0.4$。
\textbf{8.} $C\mathbf{w}_{0} = (1, 0.6)$，$\mathbf{w}_{1} = (1.10,
0.06)$，与 $(1,1)$ 的夹角：$45^{\circ} - 3.1^{\circ} = 41.9^{\circ}$。
$\mathbf{w}_{2} = (1.21, 0.13)$：$38.8^{\circ}$。$\mathbf{w}_{3} = (1.34,
0.22)$：$35.8^{\circ}$。
\textbf{9.} $w_{2}/w_{1} \to 1$，而 $|\mathbf{w}|$ 无界增长
（每步增长因子为 $1.16$）。
\textbf{10.} $C\mathbf{w} = (1.28, 1.28)$；$\mathbf{w}^{\mathsf{T}}C
\mathbf{w} = 2.05$；$\Delta\mathbf{w} = 0.1\times(1.28 - 2.05\times 0.8)
= -0.036$ 各：$\mathbf{w} = (0.76, 0.76)$，向单位
向量 $(0.71, 0.71)$ 收缩。
\textbf{11.} 联合输入 $y = 1.42$；仅输入 2 时 $y = 0.71$，减半。
\textbf{12.} 当强输入驱动细胞时，弱输入也放电，
因而与输出相关，赫布法则使之增强——
神经元学会了这种共现。
\textbf{13.} $1 - 0.85^{n}$：$0.56$，$0.80$，$0.96$。
\textbf{14.} $0.85^{n} \le 0.05$：$n \ge 18.5$，故需 $19$ 次试验。
\textbf{15.} 否：经 $n$ 次试验后达到的渐近值比例
与 $\lambda$ 无关，渐近值加倍。较大的奖励可能
提高 $\alpha$（显著性）；按\emph{绝对值}设定的行为阈值
则更早达到。
\textbf{16.} $V_{B} = 0$：被阻断。复合刺激被完全预测，
误差为零，\omterm{prop:b3:learning-memory:rescorla}{多巴胺}神经元保持沉默。
\textbf{17.} 单独的 B 不能预测任何结果，因此误差为 $\lambda$，
\omterm{prop:b3:learning-memory:rescorla}{多巴胺}爆发；$V_{B}$ 随后如习得过程那样上升，为 $\lambda(1 -
0.85^{n})$。
\textbf{18.} 药物引发的爆发即意味着“优于
预期”，作用于先行的任何事件之后；该规则每次都增加
这些线索与动作的价值，使它们获得与任何真实结果无关的价值，
觅求行为便成为强迫性的。
\textbf{19.} $0.14\times 3\times 10^{5} \approx 4.2\times 10^{4}$。
\textbf{20.} $4.2\times 10^{4}/20 = 2100$ 天，约六年；
\omterm{def:b3:learning-memory:kinds}{巩固}须将记忆\omterm{def:b3:cancer-biology:hallmarks}{转移}至皮层，遗忘则须
释放存储空间。
\textbf{21.} 若将新模式迅速写入皮层，就会覆盖
编码旧模式的突触（灾难性干扰）；缓慢而交错的
\omterm{def:b3:learning-memory:hippocampus}{重放}使皮层能将新旧信息整合，同时\omterm{def:b3:learning-memory:kinds}{海马}暂存新
记忆——一个快速缓冲
向一个慢速存储库供给。
\textbf{22.} 每个位置对应 $0.03\times 3\times 10^{5} = 9000$ 个细胞。稀疏
模式重叠很少，因而在彼此干扰损害检索之前可以存储多得多的模式；
容量随模式变得稀疏而增长快于 $N$，
模式更稀疏，但每个模式需占用更多细胞，
即以专一性相应地付出代价。
\textbf{23.} 压缩比为 $20$：运行中相隔 \qty{400}{ms} 放电的细胞，
在\omterm{def:b3:learning-memory:hippocampus}{重放}中相隔 \qty{20}{ms} 放电——处于峰电位时序
窗口之内。压缩的目的，是使实时中过慢而无法
关联的序列进入赫布法则的作用范围。
\textbf{24.} 一天时，记忆仍依赖\omterm{def:b3:learning-memory:kinds}{海马}；一个月时，
记忆已\omterm{def:b3:learning-memory:kinds}{巩固}至皮层，不再需要\omterm{def:b3:learning-memory:kinds}{海马}——
恰如 H.M. 的时间梯度性\omterm{def:b3:dna-repair:lesions}{损伤}，近期记忆已失，
而远期记忆则得以保全。
\textbf{25.} 十次开放之后，游离钙约为 \qty{4.2}{\micro mol/L}；
达到 \qty{95}{\%} 需 $19$ 次试验；大鼠 \omterm{def:b3:learning-memory:hippocampus}{CA3} 区中约有
$4\times 10^{4}$ 个模式。
\end{solution}
